\section{Introduction}

    Current applications of learning-based techniques in nuclear engineering indicate that \gls{AI} will soon undertake the role of reactor designers in optimizing the designs of current and future reactors. \gls{AI} methods, more specifically \gls{ML}, have already become important to most sub-fields of nuclear science and engineering (e.g., in-core fuel management, reactor safety, flow regime prediction, and radiation protection) to achieve an optimal core loading map, to make nuclear reactors safer, to quickly comprehend flow data, and to accurately identify critical isotopes in nuclear materials. In this context, recently published two review papers \cite{Gomez2020, Nissan2019} showed the potential, competence, and importance of these methods by exploring the recent studies on the use of learning-based methods in nuclear science and relevant fields. The articles also probed the suitability of \DIFaddbegin \DIFadd{the }\DIFaddend commonly used algorithms and algorithm selection criteria. However, \gls{AI}/\gls{ML} has not contemplated as an option to seek \DIFdelbegin \DIFdel{optimal design features }\DIFdelend \DIFaddbegin \DIFadd{to find optimal designs }\DIFaddend of nuclear reactor cores despite being a proven technique for optimization. Instead, \DIFdelbegin \DIFdel{the methods
widely employed for reactor core design optimization are }\DIFdelend iterative approaches such as individual parameter searches and \DIFdelbegin %DIFDELCMD < \gls{GA}%%%
\DIFdelend \DIFaddbegin \DIFadd{evolutionary algorithms have been widely employed for reactor core design optimization}\DIFaddend .

    Among \DIFdelbegin \DIFdel{the }\DIFdelend parametric studies recently reported, Wei \emph{et al.} \cite{Wei2018} investigated several critical feature parameters of a \gls{MSR} channel by parameter interpolation using KENO-VI in SCALE 6.1 \cite{scale2020}. In the study, \DIFdelbegin \DIFdel{the }\DIFdelend critical feature parameters are the fuel \DIFdelbegin \DIFdel{material}\DIFdelend \DIFaddbegin \DIFadd{type}\DIFaddend , $^{6}$Li concentration, composition of fuel \DIFdelbegin \DIFdel{material}\DIFdelend \DIFaddbegin \DIFadd{type}\DIFaddend , channel geometry, moderator \DIFdelbegin \DIFdel{,
moderator density}\DIFdelend \DIFaddbegin \DIFadd{and its density, }\DIFaddend and salt operation temperature. The research showed that $^{6}$Li has a strong effect on \DIFdelbegin \DIFdel{the }\DIFdelend criticality. In all different shapes of fuel channels, Wei \emph{et al.} \cite{Wei2018} observed that the turning point from under-moderation to over-moderation for the multiplication factor ($k$) is at a pitch to channel ratio ($P/D$) between 2 and 4, and a fuel volume fraction between 0 and 0.1. Increasing the amount of moderator resulted in an increase in $k$, but then further increase caused a decrease. Meanwhile, the temperature coefficient changes from negative to positive for pitch $>$ 12 cm.

    Anderson \emph{et al.} \cite{Anderson2019} assessed the importance of several parameters (fissile enrichment, lattice pitch, salt-moderator ratio\DIFdelbegin \DIFdel{and }\DIFdelend \DIFaddbegin \DIFadd{, and the }\DIFaddend number of fuel channels) in \DIFdelbegin \DIFdel{the }\DIFdelend \DIFaddbegin \DIFadd{a }\DIFaddend \gls{MSR} core by looking into the change of \DIFdelbegin \DIFdel{the
}\DIFdelend fission and removal transition matrices. In that work, the results showed that some of these parameters have an increasing impact on the average relative difference of transition matrices while \DIFdelbegin \DIFdel{the some }\DIFdelend \DIFaddbegin \DIFadd{others }\DIFaddend have a decreasing effect.

    The common deficiency of both studies is the neglect of strong interdependence of \DIFdelbegin \DIFdel{the }\DIFdelend \DIFaddbegin \DIFadd{those }\DIFaddend investigated parameters, by carrying out only \DIFaddbegin \DIFadd{a }\DIFaddend single parametric study. The effect of a single parameter on performance metrics (multiplication factor, conversion ratio, reactivity feedback coefficient, etc.) can be either beneficial or detrimental. As for more than one parameter, correlations between parameters can be seen as \DIFdelbegin \DIFdel{super positioning}\DIFdelend \DIFaddbegin \DIFadd{superpositioning}\DIFaddend , canceling out another, or acting on each other up to some extent. Correspondingly, such multi-parameter optimizations would require more sophisticated searching techniques like \DIFdelbegin %DIFDELCMD < \gls{GA}%%%
\DIFdel{, Particle Swarm, and }%DIFDELCMD < \gls{ANN} %%%
\DIFdelend \DIFaddbegin \DIFadd{genetic algorithm, particle swarm, and simulated annealing }\DIFaddend methods. Unlike parametric \DIFdelbegin \DIFdel{study}\DIFdelend \DIFaddbegin \DIFadd{studies}\DIFaddend , these sophisticated optimization techniques are very useful in finding optimal parameters, particularly for multi-objective problems\DIFaddbegin \DIFadd{, }\DIFaddend and are a common and effective way of determining optimal core loading maps \cite{Nissan2019}.

    In \DIFaddbegin \DIFadd{the }\DIFaddend case of using evolutionary algorithms in core designs, Kumar and Tsvetkov \cite{Kumar2015} used a \DIFdelbegin \DIFdel{Genetic Algorithm }\DIFdelend \DIFaddbegin \gls{GA} \DIFaddend method coupled to \DIFdelbegin \DIFdel{Regression Splines
}\DIFdelend \DIFaddbegin \DIFadd{regression splines }\DIFaddend to optimize several performance metrics such as multiplication factor, fast fission factor, thermal efficiency, and \DIFdelbegin \DIFdel{burnup }\DIFdelend \DIFaddbegin \DIFadd{burn up }\DIFaddend in the design of \DIFdelbegin \DIFdel{the }\DIFdelend Gas Cooled Fast Breeder reactor. \DIFdelbegin \DIFdel{This }\DIFdelend \DIFaddbegin \DIFadd{That }\DIFaddend method interpolates intermediate values of \DIFdelbegin \DIFdel{the
}\DIFdelend parameters and metrics. \DIFdelbegin \DIFdel{For }\DIFdelend \DIFaddbegin \DIFadd{Regarding }\DIFaddend a fuel pin cell, \DIFdelbegin \DIFdel{the optimal values of the }\DIFdelend \DIFaddbegin \DIFadd{optimal design values of }\DIFaddend radius, enrichment fraction, mass flow rate, and coolant temperature at the inlet were found to be 0.22 cm, 18.5 wt. \%, 63.5  $^o$C, and 35.9 kg/s, respectively. Zeng \emph{et al.} \cite{Zeng2020} performed a multi-objective optimization to the Advanced Burner Test Reactor (ABTR) core, a sodium-cooled fast test reactor (SFR), to get \DIFdelbegin \DIFdel{a final }\DIFdelend \DIFaddbegin \DIFadd{an ultimate }\DIFaddend core design. In \DIFdelbegin \DIFdel{the }\DIFdelend \DIFaddbegin \DIFadd{that }\DIFaddend work, an integrated version of the US Department of Energy's Nuclear Energy Advanced Modeling and Simulation (NEAMS) Workbench \cite{neams2019} along with \gls{GA} was used for core modeling and optimization. Although they found about 3150 optimal core designs within Pareto front solutions by optimizing \DIFdelbegin \DIFdel{the }\DIFdelend reactivity swing, Pu mass feed, core volume, core power, and peak fast flux, only six candidate designs passed their performance criteria based on \DIFdelbegin \DIFdel{the }\DIFdelend \DIFaddbegin \DIFadd{their }\DIFaddend defined constraints.

    \DIFaddbegin \DIFadd{Two other important studies \mbox{%DIFAUXCMD
\cite{pereira1999, pereira2003} }\hspace{0pt}%DIFAUXCMD
investigated a hypothetical reactor design in a cylindrical fuel rod geometry (i.e., composed of a moderator, clad, and fuel regions) to validate the }\gls{GA} \DIFadd{applicability and to make a comparison with classical linear optimization methods. In the first study, only two continuous feature parameters were considered for optimization. Optimal designs were completed using 3500 reactor simulations (100 populations and 35 generations). Later, a more complex version of the previous problem was repeated with continuous and discrete variables, very identical to our problem from the aspects of types and number of feature parameters. Optimization results were obtained with 150000 reactor simulations (300 populations and 500 generations).
}

    \DIFaddend Recently published studies \cite{Kim2018, Kim2020} reported the results of \DIFdelbegin \DIFdel{the
}\DIFdelend \DIFaddbegin \DIFadd{an }\DIFaddend assembly and core design at the Kyoto University Critical Assembly \DIFdelbegin \DIFdel{by means of
}\DIFdelend \DIFaddbegin \DIFadd{through }\DIFaddend an \gls{ANN} coupled to MCNP6 \cite{mcnp2017}. These studies are similar to fuel loading pattern optimization studies in terms of \DIFaddbegin \DIFadd{the }\DIFaddend optimization approach. In these works, \DIFdelbegin \DIFdel{various }\DIFdelend \DIFaddbegin \DIFadd{different }\DIFaddend core materials, \DIFdelbegin \DIFdel{numbers }\DIFdelend \DIFaddbegin \DIFadd{the number }\DIFaddend of fuel regions, and fuel assembly types were defined as vital feature parameters whereas fast flux and multiplication factor were the target metrics to be optimized.

    Regarding the use of \gls{ML} methods in depletion prediction of reactor applications, Bae \emph{et al.} \cite{Bae2020} successfully predicted Pressurized Water Reactor (PWR) \gls{UNF} composition by training the dense neural network in the Keras \cite{keras2015}. They showed that their proposed model outperforms the average recipe method, one of the transmutation methods which supplies neutronic calculations externally to the nuclear fuel cycle simulators, by yielding less than 1\% error for the \gls{UNF} inventory decay heat and activity and less than 5\% error for the important isotopes.

    As discussed previously, parametric studies which handle all feature variables independently ignore \DIFaddbegin \DIFadd{the }\DIFaddend complex interdependence of these variables and evaluate performance metrics \DIFdelbegin \DIFdel{of a design }\DIFdelend by changing only one \DIFdelbegin \DIFdel{of the feature variables }\DIFdelend \DIFaddbegin \DIFadd{feature variable }\DIFaddend at a time, thus fixing others\DIFdelbegin \DIFdel{whereas }\DIFdelend \DIFaddbegin \DIFadd{. Meanwhile, }\DIFaddend evolutionary algorithms suffer from the choices of optimizer \DIFdelbegin \DIFdel{hyper-parameters }\DIFdelend \DIFaddbegin \DIFadd{hyperparameters }\DIFaddend such as population size, number of generations\DIFaddbegin \DIFadd{, }\DIFaddend and cross-over rate \DIFdelbegin \DIFdel{, }\DIFdelend since it is unlikely to determine exactly what these parameters would be without making a preliminary calculation. \DIFaddbegin \DIFadd{Additionally, these optimization techniques do not provide any flexibility to the designers, such as making optimization without running reactor simulator again when constraints are changed. }\DIFaddend This hence results in excessive usage of computational resources. Furthermore, optimization methods \DIFdelbegin \DIFdel{frequently }\DIFdelend \DIFaddbegin \DIFadd{occasionally }\DIFaddend converge to local instead of global optima fail to achieve ideal designs that are eliminated by the optimizer's elimination and/or selection criteria. At this point, \DIFaddbegin \DIFadd{if sufficient data is supplied, }\DIFaddend the use of learning-based methods seems rational \DIFdelbegin \DIFdel{for the design of a }\DIFdelend \DIFaddbegin \DIFadd{to design a nuclear }\DIFaddend reactor because of their exceptional ability to \DIFdelbegin \DIFdel{fully scan every inch of the parameter space without missing the global optimum}\DIFdelend \DIFaddbegin \DIFadd{accurately predict performance metrics and appropriate to use with optimization tools as a reactor simulator}\DIFaddend .

    The \gls{ML} methods, which have not been considered for the sole design of a reactor core so far, will be the best option in this regard. \DIFdelbegin \DIFdel{Accordingly, we
recommend using }\DIFdelend \DIFaddbegin \DIFadd{This article compares various }\DIFaddend \gls{ML} methods \DIFdelbegin \DIFdel{to estimate the optimal values of reactor
design parameters by predicting performance metrics without running neutron transport calculations}\DIFdelend \DIFaddbegin \DIFadd{in terms of performance scores and recommends employing these methods in place of neutron transport codes for the prediction of performance metrics}\DIFaddend . This research \DIFdelbegin \DIFdel{demonstrate }\DIFdelend \DIFaddbegin \DIFadd{demonstrates }\DIFaddend that the proposed method predicts performance metrics more accurately and quickly than \DIFdelbegin \DIFdel{current
}\DIFdelend \DIFaddbegin \DIFadd{the existing }\DIFaddend optimization techniques.
